\makeatletter
\def\input@path{{preprints-template/preprints_template_latex/}}
\let\preprint@bibliographystyle\bibliographystyle
\renewcommand{\bibliographystyle}[1]{%
  \preprint@bibliographystyle{preprints-template/preprints_template_latex/#1}}
\makeatother

\documentclass[preprints,article,submit,oneauthor]{Definitions/mdpi}

\makeatletter
\let\bibliographystyle\preprint@bibliographystyle
\makeatother

\graphicspath{{preprints-template/preprints_template_latex/}}
\setlength{\headheight}{18pt}
\addtolength{\topmargin}{-6pt}

\firstpage{1}
\makeatletter
\setcounter{page}{\@firstpage}
\makeatother
\pubvolume{1}
\issuenum{1}
\articlenumber{0}
\pubyear{2026}
\copyrightyear{2026}
\datereceived{}
\daterevised{}
\dateaccepted{}
\datepublished{}

\newcounter{manuscriptresult}[section]
\renewcommand{\themanuscriptresult}{\thesection.\arabic{manuscriptresult}}
\theoremstyle{mdpi}
\newtheorem{theorem}[manuscriptresult]{Theorem}
\newtheorem{lemma}[manuscriptresult]{Lemma}
\newtheorem{proposition}[manuscriptresult]{Proposition}
\newtheorem{corollary}[manuscriptresult]{Corollary}

\theoremstyle{definition}
\newtheorem{remark}[manuscriptresult]{Remark}

\newcommand{\C}{\mathbb C}
\newcommand{\D}{\mathbb D}
\newcommand{\Mn}{\mathbb M_n(\C)}
\newcommand{\alg}{\operatorname{alg}}
\newcommand{\RePart}{\operatorname{Re}}
\newcommand{\diag}{\operatorname{diag}}
\newcommand{\iu}{\mathrm{i}}

\Title{The scalar spectral constant of the quantum annulus}
\Author{Shanmu Jin}
\AuthorNames{Shanmu Jin}
\address{Department of Neurosurgery, Peking Union Medical College
  Hospital, No.~1 Shuaifuyuan, Dongcheng District, Beijing 100730,
  China\\
jinshanmu1129@pku.edu.cn}

\abstract{
Fix $R>1$.  The quantum annulus consists of the invertible Hilbert-space
operators $A$ such that
$\lVert A\rVert,\lVert A^{-1}\rVert\leq R$.  We prove that its scalar
spectral constant for the closed annulus
$\{z\in\C:R^{-1}\leq\lvert z\rvert\leq R\}$ is exactly $2$.
The upper bound is first established for matrices.  A shifted annular
double-layer density gives a unital positive map of total mass $2$.
Applied to the Cayley transforms of a bounded scalar function, this map
produces a matrix-valued positive-real completion whose defect lies in
the adjoint algebra of an auxiliary simple-spectrum matrix.  A direct
sampling argument for the matrix Herglotz kernel then gives the sharp
norm bound.  Approximation removes the simple-spectrum and strictness
assumptions.  For arbitrary Hilbert spaces, a quantum-annulus boundary
dilation is combined with residual finite dimensionality of the full
unital free product $C(\mathbb T)*_{\C}\C^2$.  Finite cyclic weighted
shifts give the matching lower bound.  The result concerns scalar
functions; no assertion about the complete spectral constant is made.
}

\keyword{quantum annulus; spectral set; double-layer potential;
matrix-valued Herglotz kernel; residually finite-dimensional
C-star algebra}
\MSC{Primary 47A25; Secondary 47A20, 46L05}

\makeatletter
\def\maketitle{%
  \hypersetup{
    pdftitle={\@Title},
    pdfsubject={The scalar spectral constant of the quantum annulus is
      exactly two. The proof combines a positive annular double-layer
      completion, a boundary dilation, residual finite dimensionality,
      and finite cyclic weighted shifts.},
    pdfkeywords={quantum annulus; spectral set; double-layer potential;
      matrix-valued Herglotz kernel; residually finite-dimensional
      C-star algebra},
    pdfauthor={\@AuthorNames}
  }%
  \org@maketitle
}
\makeatother

\begin{document}

\section{Introduction}
\label{sec:introduction}

For $R>1$, write
\begin{equation}
  \mathbb A_R=\{z\in\C:R^{-1}<\lvert z\rvert<R\},
  \qquad
  \overline{\mathbb A}_R
  =\{z\in\C:R^{-1}\leq\lvert z\rvert\leq R\}.
  \label{eq:annulus}
\end{equation}
For a complex Hilbert space $\mathcal H$, the corresponding quantum
annulus is
\begin{equation}
  \mathbb Q\mathbb A_R(\mathcal H)
  =\left\{A\in\mathcal B(\mathcal H):
  A\text{ is invertible},\;
  \lVert A\rVert\leq R,\;
  \lVert A^{-1}\rVert\leq R\right\}.
  \label{eq:quantum-annulus}
\end{equation}
Every $A$ in this class has spectrum in
$\overline{\mathbb A}_R$.  Let $K(R)$ be the least number such that
\begin{equation}
  \lVert f(A)\rVert
  \leq K(R)\max_{z\in\overline{\mathbb A}_R}\lvert f(z)\rvert
  \label{eq:spectral-constant}
\end{equation}
for every Hilbert space $\mathcal H$, every
$A\in\mathbb Q\mathbb A_R(\mathcal H)$, and every rational function
$f$ with no pole on $\overline{\mathbb A}_R$.

Crouzeix and Greenbaum proved the dimension-free upper bound
$K(R)\leq1+\sqrt2$ by an annular double-layer argument
\cite[Section~5]{CrouzeixGreenbaum2019}.  Tsikalas proved the lower
bound $K(R)\geq2$ for every $R>1$\cite[Theorem~1.1]{Tsikalas2022}.
The annular double-layer map and its associated operator classes were
subsequently analyzed by Jury and Tsikalas
\cite{JuryTsikalas2024}.  Our main result closes the gap.

\begin{theorem}[Main theorem]
\label{thm:main}
For every $R>1$,
\begin{equation}
  K(R)=2.
  \label{eq:constant-two}
\end{equation}
Equivalently, if $\mathcal H$ is a complex Hilbert space,
$A\in\mathbb Q\mathbb A_R(\mathcal H)$, and $f$ is rational with no
pole on $\overline{\mathbb A}_R$, then
\begin{equation}
  \lVert f(A)\rVert
  \leq2\max_{z\in\overline{\mathbb A}_R}\lvert f(z)\rvert,
  \label{eq:main-estimate}
\end{equation}
and the factor $2$ cannot be decreased.
\end{theorem}

The proof has two upper-bound stages.  In finite dimensions, positivity
of a shifted annular double-layer density supplies the positive-real
completion needed in Theorem~\ref{thm:completion}.  For general Hilbert
spaces, the boundary dilation of McCullough and Pascoe
\cite[Theorem~1.1(c)]{McCulloughPascoe2023} reduces Laurent polynomials
to boundary operators; the free-product theorem of Exel and Loring
\cite[Theorem~3.2]{ExelLoring1992} then reduces their norms to finite
dimensions.  Section~\ref{sec:sharpness} gives a finite-dimensional
lower-bound sequence.

We use $I$ for the identity operator of the size dictated by context,
$X^*$ for the adjoint, and
$\RePart X=(X+X^*)/2$.  For Hermitian matrices $X,Y$, the notation
$X\preceq Y$ means that $Y-X$ is positive semidefinite.  The unital
algebra generated by an operator $X$ is denoted by $\alg(X)$.

\section{A positive-real completion theorem}
\label{sec:completion}

We begin with the finite-dimensional mechanism that produces the
constant $2$.

\begin{lemma}[Matrix Herglotz kernel]
\label{lem:herglotz-kernel}
Let $F\colon\D\to\Mn$ be analytic and suppose that
$\RePart F(w)\succeq0$ for $w\in\D$.  Then
\begin{equation}
  \mathcal L_F(w,z)
  =\frac{F(w)+F(z)^*}{1-w\overline z}
  \label{eq:Herglotz-kernel}
\end{equation}
is a positive kernel: for arbitrary
$w_0,\ldots,w_m\in\D$ and $\xi_0,\ldots,\xi_m\in\C^n$,
\begin{equation}
  \sum_{i,j=0}^m
  \xi_i^*\mathcal L_F(w_i,w_j)\xi_j\geq0.
  \label{eq:kernel-positive}
\end{equation}
\end{lemma}

\begin{proof}
Set
\[
  \eta(\zeta)=\sum_{j=0}^m
  \frac{\xi_j}{1-\overline{w_j}\zeta},
  \qquad \lvert\zeta\rvert=1.
\]
For $0\leq r<1$, entrywise integration of the power series of $F$
and of the two scalar Cauchy kernels gives, for each $i,j$,
\begin{align*}
  &\frac1{2\pi}\int_0^{2\pi}
  \frac{F(re^{\iu t})}
  {(1-w_i e^{-\iu t})(1-\overline{w_j}e^{\iu t})}\,dt
  =\frac{F(rw_i)}{1-w_i\overline{w_j}},\\
  &\frac1{2\pi}\int_0^{2\pi}
  \frac{F(re^{\iu t})^*}
  {(1-w_i e^{-\iu t})(1-\overline{w_j}e^{\iu t})}\,dt
  =\frac{F(rw_j)^*}{1-w_i\overline{w_j}}.
\end{align*}
After summing against the vectors $\xi_i,\xi_j$, these identities yield
\begin{align}
  &\frac1{2\pi}\int_0^{2\pi}
  \eta(e^{\iu t})^*\,2\RePart F(re^{\iu t})\,
  \eta(e^{\iu t})\,dt \notag\\
  &\hspace{28mm}=
  \sum_{i,j=0}^m\xi_i^*
  \frac{F(rw_i)+F(rw_j)^*}{1-w_i\overline{w_j}}\xi_j.
  \label{eq:Herglotz-integral}
\end{align}
The left side is nonnegative.  Letting $r$ increase to $1$ proves
\eqref{eq:kernel-positive}.
\end{proof}

\begin{theorem}[Positive-real completion]
\label{thm:completion}
Suppose
\begin{equation}
  B=S\diag(\beta_1,\ldots,\beta_n)S^{-1},
  \qquad
  T=S\Lambda S^{-1},
  \qquad
  \Lambda=\diag(\lambda_1,\ldots,\lambda_n),
  \label{eq:simultaneous-diagonalization}
\end{equation}
where $S$ is invertible, the numbers $\beta_1,\ldots,\beta_n$ are
distinct, and $\lvert\lambda_i\rvert\leq1$.  The numbers $\lambda_i$
need not be distinct.  Let $H\colon\D\to\Mn$ be analytic and satisfy
\begin{equation}
  H(0)=I,
  \qquad
  \RePart H(w)\succeq0\quad(w\in\D),
  \label{eq:completion-positive}
\end{equation}
and
\begin{equation}
  H(w)-(I-wT)^{-1}\in\alg(B^*)
  \quad(w\in\D).
  \label{eq:completion-algebra}
\end{equation}
Then $\lVert T\rVert\leq2$.
\end{theorem}

\begin{proof}
Put $\Gamma=S^*S\succ0$.  Since the eigenvalues of $B$ are distinct,
polynomial interpolation gives
\begin{equation}
  \alg(B^*)
  =\{(S^{-1})^*\Delta S^*: \Delta\text{ is diagonal}\}.
  \label{eq:adjoint-algebra}
\end{equation}
Consequently, there is a diagonal analytic matrix
$\Theta(w)=\diag(\theta_1(w),\ldots,\theta_n(w))$, with
$\Theta(0)=0$, such that
\begin{equation}
  \widetilde H(w):=S^*H(w)S
  =\Gamma(I-w\Lambda)^{-1}+\Theta(w)\Gamma.
  \label{eq:Htilde}
\end{equation}
The real part of $\widetilde H$ is positive, so
Lemma~\ref{lem:herglotz-kernel} applies to $\widetilde H$.

Define Hermitian matrices $P,Q,Y$ by
\begin{equation}
  P_{ij}=\frac{\Gamma_{ij}}
  {1-\overline{\lambda_i}\lambda_j/4},
  \qquad
  Q_{ij}=\frac{\Gamma_{ij}}
  {1-\overline{\lambda_i}\lambda_j/2},
  \qquad
  Y=Q-P.
  \label{eq:PQY}
\end{equation}
Fix $x=(x_1,\ldots,x_n)^T\in\C^n$.  In the kernel inequality, use
\begin{equation}
  w_i=\frac{\overline{\lambda_i}}2,
  \qquad
  \xi_i=x_i e_i
  \quad(1\leq i\leq n),
  \label{eq:sampling-points}
\end{equation}
and add the sample $w_0=0$ with vector
\begin{equation}
  \xi_0=v=-\Gamma^{-1}Px.
  \label{eq:origin-vector}
\end{equation}
Repeated sample points cause no difficulty.

The contribution of the term $\Theta(w)\Gamma$ in
\eqref{eq:Htilde} to the sampled kernel quadratic form equals
\begin{equation}
  2\RePart\sum_{i=1}^n
  \overline{x_i}\theta_i(w_i)
  \left((\Gamma v)_i+
  \sum_{j=1}^n
  \frac{\Gamma_{ij}x_j}{1-w_i\overline{w_j}}\right).
  \label{eq:theta-contribution}
\end{equation}
The expression in parentheses is
$(\Gamma v+Px)_i=0$, so the unknown term cancels exactly.

For the known term $\Gamma(I-w\Lambda)^{-1}$, compression of the
kernel block between $w_i$ and $w_j$ by $e_i$ and $e_j$ gives
\begin{equation}
  \frac{2\Gamma_{ij}}
  {(1-\overline{\lambda_i}\lambda_j/2)
   (1-\overline{\lambda_i}\lambda_j/4)}
  =(4Q-2P)_{ij}.
  \label{eq:main-sampled-block}
\end{equation}
Here the equality on the right follows from
\[
  \frac{2}{(1-a/2)(1-a/4)}
  =\frac4{1-a/2}-\frac2{1-a/4}.
\]
Likewise, direct substitution of $w_0=0$ gives
\begin{equation}
  \left[e_i^*\left(
  \Gamma(I-w_i\Lambda)^{-1}+\Gamma\right)e_j\right]_{i,j}
  =\Gamma+Q,
  \qquad
  \Gamma(I-0\Lambda)^{-1}+\Gamma=2\Gamma
  \label{eq:origin-blocks}
\end{equation}
for the known term.  Kernel positivity, evaluated on the graph
$v=-\Gamma^{-1}Px$, therefore gives
\begin{align}
  0&\leq
  \begin{pmatrix}x\\v\end{pmatrix}^{\!*}
  \begin{pmatrix}
    4Q-2P&\Gamma+Q\\
    \Gamma+Q&2\Gamma
  \end{pmatrix}
  \begin{pmatrix}x\\v\end{pmatrix} \notag\\
  &=x^*\bigl(4Y-Y\Gamma^{-1}P-P\Gamma^{-1}Y\bigr)x.
  \label{eq:Y-inequality-form}
\end{align}
Since $x$ is arbitrary,
\begin{equation}
  4Y-Y\Gamma^{-1}P-P\Gamma^{-1}Y\succeq0.
  \label{eq:Y-inequality}
\end{equation}

Balance the diagonalization by setting
\begin{equation}
  \widetilde T=\Gamma^{1/2}\Lambda\Gamma^{-1/2},
  \qquad
  \widehat P=\Gamma^{-1/2}P\Gamma^{-1/2},
  \qquad
  \widehat Y=\Gamma^{-1/2}Y\Gamma^{-1/2}.
  \label{eq:balanced-objects}
\end{equation}
Expanding the scalar denominators in \eqref{eq:PQY} gives the
norm-convergent series
\begin{align}
  \widehat P
  &=\sum_{k=0}^{\infty}4^{-k}\widetilde T^{*k}\widetilde T^k,
  \label{eq:P-series}\\
  \widehat Y
  &=\sum_{k=1}^{\infty}(2^{-k}-4^{-k})
  \widetilde T^{*k}\widetilde T^k\succeq0.
  \label{eq:Y-series}
\end{align}
Indeed, the sequence $(\widetilde T^k)$ is bounded because
$\widetilde T^k=\Gamma^{1/2}\Lambda^k\Gamma^{-1/2}$.
Congruence of \eqref{eq:Y-inequality} by $\Gamma^{-1/2}$ yields
\begin{equation}
  4\widehat Y-\widehat Y\widehat P-
  \widehat P\widehat Y\succeq0.
  \label{eq:balanced-inequality}
\end{equation}

Let $e$ be an eigenvector of $\widehat P$, with
$\widehat Pe=\alpha e$.  If $\alpha>2$, then
\eqref{eq:balanced-inequality} gives
\[
  0\leq2(2-\alpha)e^*\widehat Ye.
\]
Thus $e^*\widehat Ye=0$.  The $k=1$ term of
\eqref{eq:Y-series} is $\widetilde T^*\widetilde T/4$, so
$\widetilde Te=0$.  Equation
\eqref{eq:P-series} then gives $\widehat Pe=e$, contradicting
$\alpha>2$.  Hence
\begin{equation}
  I\preceq\widehat P\preceq2I.
  \label{eq:P-bounds}
\end{equation}
The $k=1$ term of \eqref{eq:P-series} now gives
\[
  \frac14\widetilde T^*\widetilde T\preceq\widehat P-I\preceq I,
\]
and therefore $\lVert\widetilde T\rVert\leq2$.  Finally, the polar
decomposition $S=U_S\Gamma^{1/2}$ gives
$T=U_S\widetilde TU_S^*$, so
$\lVert T\rVert\leq2$.
\end{proof}

\section{The annular double-layer completion}
\label{sec:annular-completion}

Let $\mathcal A(\mathbb A_R)$ denote the functions holomorphic on
$\mathbb A_R$ and continuous on $\overline{\mathbb A}_R$.  We now
construct the completion in Theorem~\ref{thm:completion}.  The shifted
density below is the annular double-layer density underlying the
contractive map in \cite[Section~5]{CrouzeixGreenbaum2019}; see also
\cite[Theorem~3.12]{JuryTsikalas2024}.  We include the factorization,
mass, and cancellation identities needed here.

Let $B\in\Mn$ satisfy
\begin{equation}
  \lVert B\rVert<R,
  \qquad
  \lVert B^{-1}\rVert<R.
  \label{eq:strict-annulus}
\end{equation}
Orient the outer circle of $\partial\mathbb A_R$ counterclockwise and
the inner circle clockwise.  On either component, let $z=z(s)$ be an
oriented arclength parametrization.  Define
\begin{align}
  \mu_B(z)
  &=\frac1{2\pi\iu}\left[
  z'(s)(zI-B)^{-1}
  -\overline{z'(s)}(\overline zI-B^*)^{-1}\right],
  \label{eq:mu-density}\\
  \kappa_R(z)
  &=\frac1{2\pi\iu}\frac{z'(s)}z,
  \qquad
  D_B(z)=\mu_B(z)-\kappa_R(z)I.
  \label{eq:shifted-density}
\end{align}
Both $\mu_B(z)$ and $D_B(z)$ are self-adjoint.

\begin{lemma}[Positive shifted density]
\label{lem:positive-density}
Under \eqref{eq:strict-annulus},
\begin{equation}
  D_B(z)\succeq0
  \quad(z\in\partial\mathbb A_R),
  \qquad
  \int_{\partial\mathbb A_R}D_B(z)\,ds=2I.
  \label{eq:density-positive-mass}
\end{equation}
Moreover, for every $h\in\mathcal A(\mathbb A_R)$,
\begin{equation}
  \int_{\partial\mathbb A_R}
  h(z)\kappa_R(z)\,ds=0.
  \label{eq:kappa-cancellation}
\end{equation}
\end{lemma}

\begin{proof}
On the outer circle, write $z=Re^{\iu t}$.  Then
$\kappa_R(z)=1/(2\pi R)$, and direct multiplication by
$zI-B$ and $\overline zI-B^*$ gives
\begin{equation}
  2\pi R D_B(z)
  =(zI-B)^{-1}(R^2I-BB^*)
  (\overline zI-B^*)^{-1}\succeq0.
  \label{eq:outer-factorization}
\end{equation}
On the inner circle, put $\rho=R^{-1}$ and use the clockwise
parametrization $z=\rho e^{-\iu t}$.  Now
$\kappa_R(z)=-1/(2\pi\rho)$ and
\begin{equation}
  2\pi\rho D_B(z)
  =(zI-B)^{-1}(BB^*-\rho^2I)
  (\overline zI-B^*)^{-1}\succeq0.
  \label{eq:inner-factorization}
\end{equation}
The middle factors are positive because
$BB^*\preceq R^2I$ and $BB^*\succeq R^{-2}I$.

The matrix Cauchy formula gives
\begin{equation}
  \frac1{2\pi\iu}\int_{\partial\mathbb A_R}
  z'(s)(zI-B)^{-1}\,ds
  =\frac1{2\pi\iu}\int_{\partial\mathbb A_R}
  (zI-B)^{-1}\,dz=I.
  \label{eq:first-resolvent-mass}
\end{equation}
Under the change of variable $\zeta=\overline z$, the conjugate of the
oriented annular boundary has the reverse orientation.  Therefore
\begin{equation}
  -\frac1{2\pi\iu}\int_{\partial\mathbb A_R}
  \overline{z'(s)}(\overline zI-B^*)^{-1}\,ds
  =-\frac1{2\pi\iu}\int_{-\partial\mathbb A_R}
  (\zeta I-B^*)^{-1}\,d\zeta=I.
  \label{eq:second-resolvent-mass}
\end{equation}
Adding \eqref{eq:first-resolvent-mass} and
\eqref{eq:second-resolvent-mass} yields
\begin{equation}
  \int_{\partial\mathbb A_R}\mu_B(z)\,ds=2I.
  \label{eq:mu-mass}
\end{equation}
The winding numbers of the outer and inner circles are $1$ and $-1$,
respectively, so
\[
  \int_{\partial\mathbb A_R}\kappa_R(z)\,ds
  =\frac1{2\pi\iu}\int_{\partial\mathbb A_R}\frac{dz}{z}=0.
\]
This proves the mass identity.  Finally,
\[
  \int_{\partial\mathbb A_R}h(z)\kappa_R(z)\,ds
  =\frac1{2\pi\iu}
  \int_{\partial\mathbb A_R}\frac{h(z)}z\,dz=0
\]
by Cauchy's theorem on the annulus.
\end{proof}

Define
\begin{equation}
  \Phi_B(\varphi)
  =\frac12\int_{\partial\mathbb A_R}
  \varphi(z)D_B(z)\,ds,
  \qquad
  \varphi\in C(\partial\mathbb A_R).
  \label{eq:positive-map}
\end{equation}
Lemma~\ref{lem:positive-density} makes $\Phi_B$ a unital positive map;
in particular, it is bounded and $*$-preserving.  For
$h\in\mathcal A(\mathbb A_R)$, define
\begin{equation}
  (\mathcal C\overline h)(\zeta)
  =\frac1{2\pi\iu}
  \int_{\partial\mathbb A_R}
  \frac{\overline{h(z)}}{z-\zeta}\,dz,
  \qquad \zeta\in\mathbb A_R.
  \label{eq:conjugate-Cauchy-transform}
\end{equation}
This is holomorphic in $\mathbb A_R$.  Splitting the two terms of
$\mu_B$ and using \eqref{eq:kappa-cancellation} gives
\begin{align*}
  \int_{\partial\mathbb A_R}h(z)D_B(z)\,ds
  &=\frac1{2\pi\iu}\int_{\partial\mathbb A_R}
  h(z)(zI-B)^{-1}\,dz\\
  &\quad-\frac1{2\pi\iu}\int_{\partial\mathbb A_R}
  h(z)(\overline zI-B^*)^{-1}\,d\overline z\\
  &=h(B)+
  \left[\frac1{2\pi\iu}\int_{\partial\mathbb A_R}
  \overline{h(z)}(zI-B)^{-1}\,dz\right]^*.
\end{align*}
The bracket is $(\mathcal C\overline h)(B)$ by
\eqref{eq:conjugate-Cauchy-transform} and the holomorphic functional
calculus.  Since the left side is $2\Phi_B(h)$, we obtain
\begin{equation}
  2\Phi_B(h)=h(B)+(\mathcal C\overline h)(B)^*.
  \label{eq:symmetrized-calculus}
\end{equation}

\begin{proposition}[Annular Cayley completion]
\label{prop:annular-completion}
Suppose that $B\in\Mn$ satisfies \eqref{eq:strict-annulus} and has
$n$ distinct eigenvalues.  If $f\in\mathcal A(\mathbb A_R)$, then
\begin{equation}
  \lVert f(B)\rVert
  \leq2\max_{z\in\overline{\mathbb A}_R}\lvert f(z)\rvert.
  \label{eq:simple-spectrum-estimate}
\end{equation}
\end{proposition}

\begin{proof}
Put
$M=\max_{z\in\overline{\mathbb A}_R}\lvert f(z)\rvert$.
If $M=0$, then $f=0$ and the result is immediate.  Replacing $f$ by
$f/M$, we may assume that $M\leq1$.
Put $T=f(B)$.  For $w\in\D$, define
\begin{equation}
  c_w(z)=\frac{1+w f(z)}{1-w f(z)},
  \qquad
  H(w)=\Phi_B(c_w).
  \label{eq:Cayley-family}
\end{equation}
The series
$c_w=1+2\sum_{m\geq1}w^m f^m$ converges locally uniformly in the
supremum norm on $\overline{\mathbb A}_R$.  Thus $H$ is analytic and
$H(0)=I$.  Since
\[
  \RePart c_w(z)
  =\frac{1-\lvert w f(z)\rvert^2}
  {\lvert1-w f(z)\rvert^2}\geq0,
\]
positivity and $*$-preservation of $\Phi_B$ give
$\RePart H(w)\succeq0$.

Equation \eqref{eq:symmetrized-calculus} gives
\begin{equation}
  2H(w)=c_w(B)+(\mathcal C\overline{c_w})(B)^*.
  \label{eq:Cayley-symmetrization}
\end{equation}
Since
\[
  c_w(B)=(I+wT)(I-wT)^{-1}
  =2(I-wT)^{-1}-I,
\]
we obtain
\begin{equation}
  H(w)-(I-wT)^{-1}
  =\frac12\left((\mathcal C\overline{c_w})(B)^*-I\right)
  \in\alg(B^*).
  \label{eq:annular-algebra-completion}
\end{equation}
Indeed, a holomorphic function of a finite matrix is a polynomial in
that matrix, by interpolation with multiplicities, and hence belongs
to its generated unital algebra.

Choose an eigenbasis
$B=S\diag(\beta_1,\ldots,\beta_n)S^{-1}$.  Then
\[
  T=S\diag(f(\beta_1),\ldots,f(\beta_n))S^{-1},
  \qquad \lvert f(\beta_i)\rvert\leq1.
\]
Theorem~\ref{thm:completion}, applied to
\eqref{eq:annular-algebra-completion}, yields
$\lVert f(B)\rVert\leq2$ under the normalization $M\leq1$.
Scaling back by $M$ proves \eqref{eq:simple-spectrum-estimate}.
\end{proof}

\section{The finite-dimensional estimate}
\label{sec:finite-dimensions}

\begin{theorem}[Finite dimensions]
\label{thm:finite-dimensional}
Let $A\in\Mn$ be invertible and satisfy
$\lVert A\rVert,\lVert A^{-1}\rVert\leq R$.  If $f$ is rational and
has no pole on $\overline{\mathbb A}_R$, then
\begin{equation}
  \lVert f(A)\rVert
  \leq2\max_{z\in\overline{\mathbb A}_R}\lvert f(z)\rvert.
  \label{eq:finite-dimensional-estimate}
\end{equation}
\end{theorem}

\begin{proof}
Write the polar decomposition as $A=U_A\lvert A\rvert$.  Because $A$
is invertible, $U_A$ is unitary and
\begin{equation}
  R^{-1}I\preceq\lvert A\rvert\preceq RI.
  \label{eq:modulus-bounds}
\end{equation}
For $0<\varepsilon<1$, put
\begin{equation}
  A_\varepsilon
  =U_A\bigl((1-\varepsilon)\lvert A\rvert+\varepsilon I\bigr).
  \label{eq:strictification}
\end{equation}
Since $R>1$, \eqref{eq:modulus-bounds} implies
\begin{equation}
  \lVert A_\varepsilon\rVert<R,
  \qquad
  \lVert A_\varepsilon^{-1}\rVert<R,
  \qquad
  A_\varepsilon\longrightarrow A.
  \label{eq:strictification-bounds}
\end{equation}
Indeed, the spectrum of the positive factor in
\eqref{eq:strictification} lies in the open interval $(R^{-1},R)$.

The strict inequalities in \eqref{eq:strictification-bounds} persist
under sufficiently small matrix perturbations.  Matrices with distinct
eigenvalues are dense: the discriminant of the characteristic
polynomial is a nonzero polynomial in the matrix entries, so its zero
set has empty interior.  Hence, for each $\varepsilon$, there are
simple-spectrum matrices $B_{\varepsilon,k}\to A_\varepsilon$ that
satisfy \eqref{eq:strict-annulus}.  Proposition~\ref{prop:annular-completion}
gives
\[
  \lVert f(B_{\varepsilon,k})\rVert
  \leq2\max_{z\in\overline{\mathbb A}_R}\lvert f(z)\rvert.
\]
First let $k\to\infty$ and then let $\varepsilon\downarrow0$.
Continuity of rational matrix evaluation proves
\eqref{eq:finite-dimensional-estimate}.
\end{proof}

\section{Passage to arbitrary Hilbert spaces}
\label{sec:all-Hilbert-spaces}

We use two external structural theorems in this section.  In our
normalization, the boundary-dilation theorem of McCullough and Pascoe
\cite[Theorem~1.1 and Proposition~3.2]{McCulloughPascoe2023} says that
every $A\in\mathbb Q\mathbb A_R(\mathcal H)$ has a Laurent dilation to
an invertible operator $J\in\mathcal B(\mathcal L)$: there is an
isometry $V\colon\mathcal H\to\mathcal L$ such that
\begin{equation}
  A^m=V^*J^mV
  \quad(m\in\mathbb Z),
  \label{eq:Laurent-dilation}
\end{equation}
and
\begin{equation}
  (R^2+R^{-2})I-J^*J-(J^*J)^{-1}=0.
  \label{eq:boundary-equation}
\end{equation}
The second input is the theorem of Exel and Loring
\cite[Theorem~3.2]{ExelLoring1992}: the full unital free product of two
unital residually finite-dimensional $C^*$-algebras is residually
finite-dimensional.

Set
\begin{equation}
  \mathfrak F=C(\mathbb T)*_{\C}\C^2.
  \label{eq:free-product}
\end{equation}
Let $u_0\in\mathfrak F$ be the canonical unitary from $C(\mathbb T)$
and let $p_0\in\mathfrak F$ be the canonical projection from $\C^2$.
Define
\begin{equation}
  j_R=u_0\bigl(Rp_0+R^{-1}(1-p_0)\bigr).
  \label{eq:universal-boundary-element}
\end{equation}
Its inverse is
\begin{equation}
  j_R^{-1}=\bigl(R^{-1}p_0+R(1-p_0)\bigr)u_0^*.
  \label{eq:universal-boundary-inverse}
\end{equation}
The algebra $C(\mathbb T)$ is residually finite-dimensional because
point evaluations separate its elements, and $\C^2$ is
finite-dimensional.  Thus $\mathfrak F$ is residually
finite-dimensional.  Equivalently,
\begin{equation}
  \lVert x\rVert
  =\sup_{\pi}\lVert\pi(x)\rVert
  \qquad(x\in\mathfrak F),
  \label{eq:RFD-norm}
\end{equation}
where $\pi$ ranges over finite-dimensional unital
$*$-representations.  Indeed, Exel and Loring's characterization
\cite[Theorem~2.4]{ExelLoring1992} makes their direct sum faithful and
hence isometric.

\begin{lemma}[Boundary reduction]
\label{lem:boundary-reduction}
If $J$ satisfies \eqref{eq:boundary-equation}, then every Laurent
polynomial $q$ satisfies
\begin{equation}
  \lVert q(J)\rVert
  \leq
  \sup_{n\geq1}\;
  \sup_{B\in\mathbb Q\mathbb A_R(\C^n)}
  \lVert q(B)\rVert.
  \label{eq:boundary-reduction}
\end{equation}
\end{lemma}

\begin{proof}
Put $Z=J^*J$.  Multiplication of \eqref{eq:boundary-equation} by $Z$
gives
\[
  (Z-R^2I)(Z-R^{-2}I)=0.
\]
The spectral theorem therefore shows that the spectrum of $Z$ is
contained in $\{R^2,R^{-2}\}$.  Hence
\begin{equation}
  Z=R^2E+R^{-2}(I-E)
  \label{eq:two-point-modulus}
\end{equation}
for a projection $E$.  The polar decomposition of $J$ is therefore
\begin{equation}
  J=U_J\bigl(RE+R^{-1}(I-E)\bigr)
  \label{eq:boundary-polar-form}
\end{equation}
with $U_J$ unitary.  The universal property of \eqref{eq:free-product}
gives a unital $*$-representation
$\rho_J\colon\mathfrak F\to\mathcal B(\mathcal L)$ such that
$\rho_J(j_R)=J$.  Thus
\[
  \lVert q(J)\rVert\leq\lVert q(j_R)\rVert.
\]
By \eqref{eq:RFD-norm}, the last norm is the supremum of
$\lVert q(\pi(j_R))\rVert$ over finite-dimensional unital
$*$-representations $\pi$.  Each matrix
\[
  \pi(j_R)=\pi(u_0)
  \bigl(R\pi(p_0)+R^{-1}(I-\pi(p_0))\bigr)
\]
has norm at most $R$.  Its inverse is
\[
  \bigl(R^{-1}\pi(p_0)+R(I-\pi(p_0))\bigr)\pi(u_0)^*,
\]
which also has norm at most $R$.  Thus $\pi(j_R)$ belongs to the
finite-dimensional quantum annulus, proving
\eqref{eq:boundary-reduction}.
\end{proof}

\begin{proof}[Proof of the upper bound in Theorem~\ref{thm:main}]
Let $A\in\mathbb Q\mathbb A_R(\mathcal H)$ and let $J,V$ be supplied
by \eqref{eq:Laurent-dilation}--\eqref{eq:boundary-equation}.  If
$q(z)=\sum_{m=-N}^N c_mz^m$ is a Laurent polynomial, then
$q(A)=V^*q(J)V$.  Lemma~\ref{lem:boundary-reduction} and
Theorem~\ref{thm:finite-dimensional} give
\begin{align}
  \lVert q(A)\rVert
  &\leq\lVert q(J)\rVert \notag\\
  &\leq
  \sup_{n\geq1}\;
  \sup_{B\in\mathbb Q\mathbb A_R(\C^n)}\lVert q(B)\rVert
  \leq2\max_{z\in\overline{\mathbb A}_R}\lvert q(z)\rvert.
  \label{eq:Laurent-estimate}
\end{align}

Now let $f$ be rational with no pole on
$\overline{\mathbb A}_R$.  There are radii
$r_0<R^{-1}<R<r_1$ such that $f$ is holomorphic on
$r_0<\lvert z\rvert<r_1$.  Its Laurent series on this annulus has
partial sums $q_k$ satisfying
\begin{equation}
  \max_{z\in\overline{\mathbb A}_R}
  \lvert q_k(z)-f(z)\rvert\longrightarrow0.
  \label{eq:Laurent-uniform}
\end{equation}
The convergence also holds on two circles surrounding
$\sigma(A)$ inside the larger annulus, so the Cauchy integral formula
for the holomorphic functional calculus gives
$q_k(A)\to f(A)$ in operator norm.  Passing to the limit in
\eqref{eq:Laurent-estimate} proves \eqref{eq:main-estimate}.
\end{proof}

\section{Sharpness}
\label{sec:sharpness}

The lower bound is approached by finite cyclic weighted shifts.  This
is a finite-dimensional version of the weighted-shift construction in
\cite[Section~2]{Tsikalas2022}.

For $n\geq1$, let $W_n$ act on $\C^{2n}$ with standard basis
$e_0,\ldots,e_{2n-1}$ by
\begin{equation}
  W_ne_k=\omega_ke_{k+1\pmod{2n}},
  \qquad
  \omega_k=
  \begin{cases}
    R,&0\leq k<n,\\
    R^{-1},&n\leq k<2n.
  \end{cases}
  \label{eq:cyclic-shift}
\end{equation}
The cyclic permutation is unitary, so the singular values of $W_n$
are the numbers $\omega_k$.  Consequently,
\begin{equation}
  \lVert W_n\rVert=\lVert W_n^{-1}\rVert=R.
  \label{eq:shift-annulus}
\end{equation}
For the Laurent polynomial
\begin{equation}
  g_n(z)=R^{-n}(z^n+z^{-n}),
  \label{eq:lower-function}
\end{equation}
the maximum-modulus principle, applied on the annulus, gives
\begin{equation}
  \max_{z\in\overline{\mathbb A}_R}\lvert g_n(z)\rvert
  =1+R^{-2n}.
  \label{eq:lower-function-norm}
\end{equation}
Indeed, on either boundary circle the triangle inequality gives this
upper bound, and equality holds at a positive real boundary point.
On the other hand, the two half-cycles in \eqref{eq:cyclic-shift} give
\[
  W_n^ne_0=W_n^{-n}e_0=R^ne_n.
\]
Hence
\begin{equation}
  \lVert g_n(W_n)\rVert\geq2.
  \label{eq:lower-operator-norm}
\end{equation}
It follows from \eqref{eq:shift-annulus}--\eqref{eq:lower-operator-norm}
that
\[
  K(R)\geq\frac{2}{1+R^{-2n}}.
\]
Letting $n\to\infty$ proves $K(R)\geq2$.  Together with the upper
bound, this completes the proof of Theorem~\ref{thm:main}.

\begin{remark}[Scalar scope]
The theorem identifies the scalar spectral constant.  It makes no
assertion about the complete spectral constant.
\end{remark}

\section*{AI-assistance disclosure}

OpenAI ChatGPT assisted with checking the kernel-sampling calculation,
the annular boundary factorizations, the source-to-claim matching, and
the exposition and typesetting of this manuscript.  The human author
assumes full responsibility for the correctness of all mathematical
statements and for the integrity and accuracy of all citations.

\reftitle{References}
\bibliography{the_quantum_annulus_is_a_2_spectral_set}

\end{document}
